\section{Finite projections and compatible rectangles}

A finite rectangle selects temporal and spatial orders directly from the
mixed jet, retaining its velocity, pressure, and prescribed-force coordinates
and their governing relationships.  The datum-generated intersection and the
critical-height rectangle construction are independent readings of that same
jet.  The first produces the whole-terminal memberships of one forced history;
restriction and monotone exhaustion of finitely many selected memberships give
the rectangle norm.  The second uses the viscous tether in the complete
equation to estimate the complete rows carried by that rectangle.  We now
derive those identities with the nonlinear, pressure, and force faces retained.


\begin{proposition}[Kinetic-energy identity]
\label{prop:kinetic-energy}
For every \(0\le t<T_*\),
\begin{equation}
 \frac12\norm{u(t)}_2^2
 +\nu\int_0^t\norm{\nabla u(s)}_2^2\dd s
 =\frac12\norm{u_0}_2^2
 +\int_0^t\ip{f(s)}{u(s)}_{L^2}\dd s.
 \label{eq:kinetic-energy}
\end{equation}
Put
\[
 K_f(T):=\norm{u_0}_2+\int_0^T\norm{\Pp f(s)}_2\dd s.
\]
Consequently, for \(0<T<T_*\),
\begin{align}
 \norm u_{L^\infty(0,T;L^2)}&\le K_f(T),
 \label{eq:energy-Linf}\\
 \norm u_{L^2(0,T;H^1)}^2
 &\le\left(T+\frac1{2\nu}\right)K_f(T)^2.
 \label{eq:energy-L2H1}
\end{align}
\end{proposition}

\begin{proof}
Take the \(L^2\) pairing of the momentum equation with \(u\).  Integration
by parts and \(\nabla\cdot u=0\) give
\[
 \ip{(u\cdot\nabla)u}{u}_{L^2}=0,
 \qquad
 \ip{\nabla p}{u}_{L^2}=0,
 \qquad
 \ip{\Delta u}{u}_{L^2}=-\norm{\nabla u}_2^2.
\]
Since \(u\) is divergence free,
\(\ip f u=\ip{\Pp f}u\).  Integration in time proves
\eqref{eq:kinetic-energy}.  Applying the identity first to
\((\norm u_2^2+\varepsilon)^{1/2}\) and then letting
\(\varepsilon\downarrow0\) gives \(\norm{u(t)}_2\le K_f(t)\).  Moreover,
\[
 \int_0^T K_f(t)\norm{\Pp f(t)}_2\dd t
 =\frac12\bigl(K_f(T)^2-\norm{u_0}_2^2\bigr),
\]
so \eqref{eq:kinetic-energy} gives
\(\int_0^T\norm{\nabla u}_2^2\dd t\le K_f(T)^2/(2\nu)\).
Adding \(\int_0^T\norm{u(s)}_2^2\dd s\) proves
\eqref{eq:energy-L2H1}.
\end{proof}

Pairing a differentiated velocity equation with a spatial derivative of
that row expresses the change of its Sobolev energy through viscous
dissipation, the full transport--pressure contribution, and the prescribed
force work. Repeated time differentiation then exposes higher spatial energies,
with every differentiated nonlinear and force term retained in each identity.

\begin{definition}[Completed height rows]
\label{def:critical-height-rows}
For \(r\in\Nzero\) and \(s\ge0\), set
\begin{align}
 E_{r,s}(t)&:=\frac12\norm{\Lambda^sV_r(t)}_2^2,
 \label{eq:temporal-row-spatial-energy}\\
 \cP_{r,s}(t)&:=-\operatorname{Re}
 \ip{\Lambda^sV_r(t)}
 {\Lambda^s\!\left(
   \sum_{a=0}^{r}\binom ra(V_a\cdot\nabla)V_{r-a}+\nabla Q_r
 \right)(t)}_{L^2}.
 \label{eq:temporal-row-pressure-complete-transfer}
 \\
 \mathcal W_{r,s}(t)&:=\operatorname{Re}
 \ip{\Lambda^sV_r(t)}{\Lambda^sF_r(t)}_{L^2}.
 \label{eq:temporal-row-force-work}
\end{align}
For the base row write \(E_s:=E_{0,s}\), \(\cP_s:=\cP_{0,s}\),
\(\mathcal W_s:=\mathcal W_{0,s}\), and \(H:=E_{1/2}\).
\end{definition}

Differentiating this quadratic height along the classical history gives
\[
 H^{(r)}(t)=\frac12\sum_{a=0}^r\binom ra
 \ip{\Lambda^{1/2}V_a(t)}{\Lambda^{1/2}V_{r-a}(t)}_{L^2}.
\]
On every strict interval the mixed jet makes this contraction an identity of
the same forced history.  The following recurrence then exposes its spatial
face without discarding the transfer or force work.  Its whole-terminal use
is made only after the datum-generated intersection has placed the
corresponding mixed rows on \((0,T_*)\); no intersection membership is used
to derive the rectangle relationship itself.

\begin{proposition}[Completed height recurrence in every temporal row]
\label{prop:completed-height-recurrence}
For every strict classical interval \((0,T)\), every \(r\in\Nzero\), and
every \(s\ge0\),
\begin{equation}
 \partial_tE_{r,s}=\cP_{r,s}+\mathcal W_{r,s}-2\nu E_{r,s+1}.
 \label{eq:sobolev-height-balance}
\end{equation}
For every \(k\in\N\),
\begin{equation}
 \partial_t^kE_{r,1/2}
 =(-2\nu)^kE_{r,k+1/2}
 +\sum_{j=0}^{k-1}(-2\nu)^j
   \partial_t^{k-1-j}
   (\cP_{r,j+1/2}+\mathcal W_{r,j+1/2}).
 \label{eq:completed-height-every-temporal-row}
\end{equation}
In the base row \(r=0\),
\begin{equation}
 H^{(k)}
 =(-2\nu)^kE_{k+1/2}
 +\sum_{j=0}^{k-1}(-2\nu)^j
   \partial_t^{k-1-j}(\cP_{j+1/2}+\mathcal W_{j+1/2}),
 \label{eq:completed-height-recurrence}
\end{equation}
and the highest spatial energy is the exact coordinate
\begin{equation}
 E_{k+1/2}
 =\frac{H^{(k)}-
   \displaystyle\sum_{j=0}^{k-1}(-2\nu)^j
   \partial_t^{k-1-j}(\cP_{j+1/2}+\mathcal W_{j+1/2})}
 {(-2\nu)^k}.
 \label{eq:completed-height-row}
\end{equation}
\end{proposition}

\begin{proof}
Pair the temporal recurrence for \(V_r\) with \(\Lambda^{2s}V_r\).
The pressure-complete nonlinear sum is exactly \(\cP_{r,s}\), the force
pairing is exactly \(\mathcal W_{r,s}\), and
\[
 \nu\operatorname{Re}\ip{\Lambda^sV_r}{\Lambda^s\Delta V_r}_{L^2}
 =-\nu\norm{\Lambda^{s+1}V_r}_2^2=-2\nu E_{r,s+1}.
\]
This proves \eqref{eq:sobolev-height-balance}.  Induction on \(k\) gives
the stronger identity
\begin{equation}
 \partial_t^kE_{r,s}
 =(-2\nu)^kE_{r,s+k}
 +\sum_{j=0}^{k-1}(-2\nu)^j
   \partial_t^{k-1-j}(\cP_{r,s+j}+\mathcal W_{r,s+j}).
 \label{eq:all-height-recurrence}
\end{equation}
Indeed, differentiation of the order-\(k\) identity followed by
\(\partial_tE_{r,s+k}=\cP_{r,s+k}+\mathcal W_{r,s+k}
-2\nu E_{r,s+k+1}\) produces the
order-\(k+1\) identity.  Taking \(s=1/2\) proves
\eqref{eq:completed-height-every-temporal-row}.  Taking \(r=0\) gives
\eqref{eq:completed-height-recurrence}; since \(\nu>0\), division by
\((-2\nu)^k\) gives \eqref{eq:completed-height-row}.
\end{proof}

\begin{definition}[Completed critical-height coordinates]
\label{def:completed-critical-height-coordinates}
Set
\begin{equation}
 R_0:=H=E_{1/2},
 \qquad
 R_k
 :=\frac{H^{(k)}-
   \displaystyle\sum_{j=0}^{k-1}(-2\nu)^j
   \partial_t^{k-1-j}(\cP_{j+1/2}+\mathcal W_{j+1/2})}
 {(-2\nu)^k}
 \quad(k\in\N).
 \label{eq:completed-critical-coordinate}
\end{equation}
Proposition~\ref{prop:completed-height-recurrence} gives the exact identity
\begin{equation}
 R_k=E_{k+1/2}
 =\frac12\norm{\Lambda^{k+1/2}u}_2^2
 \qquad(k\in\Nzero).
 \label{eq:completed-critical-coordinate-identity}
\end{equation}
\par
\end{definition}

\begin{proposition}[Complete forced rectangle cell]
\label{prop:complete-forced-rectangle-cell}
For every temporal row \(n\in\Nzero\), every \(s\ge0\), and every strict
classical interval, the projected differentiated equation is
\begin{equation}
 \Lambda^sV_{n+1}+\nu\Lambda^{s+2}V_n
 =\Lambda^s\Pp(F_n-B_n),
 \qquad
 B_n:=\sum_{a=0}^n\binom na(V_a\cdot\nabla)V_{n-a}.
 \label{eq:complete-forced-rectangle-cell}
\end{equation}
Its square is the viscous diagonal identity
\begin{equation}
 \norm{\Lambda^sV_{n+1}}_2^2
 +\nu^2\norm{\Lambda^{s+2}V_n}_2^2
 +\nu\frac{\dd}{\dd t}\norm{\Lambda^{s+1}V_n}_2^2
 =\norm{\Lambda^s\Pp(F_n-B_n)}_2^2.
 \label{eq:complete-forced-rectangle-square}
\end{equation}
Before projection, the gradient face is orthogonal to both divergence-free
faces.  Consequently,
\begin{align}
 &\norm{V_{n+1}}_{\dot H^s}^2
 +\nu^2\norm{\Delta V_n}_{\dot H^s}^2
 +\norm{\nabla Q_n}_{\dot H^s}^2
 +\nu\frac{\dd}{\dd t}\norm{\nabla V_n}_{\dot H^s}^2
 \notag\\
 &\hspace{4cm}=\norm{F_n-B_n}_{\dot H^s}^2.
 \label{eq:complete-forced-pressure-square}
\end{align}
Thus, for \(0<S<T_*\),
\begin{align}
 &\nu\norm{\Lambda^{s+1}V_n(S)}_2^2
 +\int_0^S\!\left(
   \norm{\Lambda^sV_{n+1}}_2^2
   +\nu^2\norm{\Lambda^{s+2}V_n}_2^2
   +\norm{\Lambda^s\nabla Q_n}_2^2\right)\dd t
 \notag\\
 &\qquad=
 \nu\norm{\Lambda^{s+1}V_n^0}_2^2
 +\int_0^S\norm{\Lambda^s(F_n-B_n)}_2^2\dd t.
 \label{eq:complete-forced-rectangle-integrated}
\end{align}
No term of the differentiated momentum equation has been removed or replaced
in these identities.
\end{proposition}

\begin{proof}
Apply \(\Pp\Lambda^s\) to the temporal row in
Proposition~\ref{prop:formal-mixed-recurrences}. Since
\(-\Delta=\Lambda^2\), moving the viscous term to the left gives
\eqref{eq:complete-forced-rectangle-cell} with a positive sign.
On squaring, the cross term is
\[
 2\nu\operatorname{Re}
 \ip{\Lambda^sV_{n+1}}{\Lambda^{s+2}V_n}_{L^2}
 =\nu\frac{\dd}{\dd t}\norm{\Lambda^{s+1}V_n}_2^2,
\]
which proves \eqref{eq:complete-forced-rectangle-square}.  The complementary
projection gives
\(\nabla Q_n=(I-\Pp)(F_n-B_n)\).  Orthogonality of \(\Pp\) and
\(I-\Pp\) proves \eqref{eq:complete-forced-pressure-square}; integration in
time gives \eqref{eq:complete-forced-rectangle-integrated}.
\end{proof}

\begin{definition}[Finite mixed rectangle]
\label{def:finite-rectangle}
Let \(0<T<T_*\) and \(N,M\in\Nzero\).  Define
\begin{align}
 \mathfrak R_{N,M}(u;T)
 :={}&\sum_{n=0}^{N}
 \norm{V_n}_{L^2(0,T;H^{M+1})}^2
 +\sum_{n=0}^{N}
 \norm{V_{n+1}}_{L^2(0,T;H^{M-1})}^2,
 \label{eq:finite-rectangle-norm}\\
 \mathfrak Q_{N,M}(p;T)
 :={}&\sum_{n=0}^{N}
 \left(
 \norm{Q_n}_{L^1(0,T;H^M)}
 +\norm{\nabla Q_n}_{L^1(0,T;H^{M-1})}
 \right).
 \label{eq:finite-pressure-rectangle-norm}
 \\
 \mathfrak F_{N,M}(f;T)
 :={}&\sum_{n=0}^{N}\left(
 \norm{F_n}_{L^1(0,T;H^{M-1})}
 +\mathfrak F^{\rm p}_{n,0,M}(0,T)
 \right).
 \label{eq:finite-force-rectangle-norm}
\end{align}
The pressure-complete forced record is
\[
 \cR_{N,M}[u_0,\nu,f;T]
 :=\left(
 (V_n|_{(0,T)})_{n=0}^{N},
 (V_{n+1}|_{(0,T)})_{n=0}^{N};
 (Q_n|_{(0,T)})_{n=0}^{N};
 (F_n|_{(0,T)})_{n=0}^{N}
 \right),
\]
with the pressure-gradient and force coordinates retained through
\eqref{eq:finite-pressure-rectangle-norm}--\eqref{eq:finite-force-rectangle-norm}
and with every row constrained by
Proposition~\ref{prop:formal-mixed-recurrences}.
\end{definition}

\Needspace{14\baselineskip}
\begin{theorem}[Datum-generated whole-terminal rectangles and complete-equation estimates]
\label{thm:datum-rectangles}

Let \(\nu>0\), \(u_0\in\HsSigma(\R^3)\), prescribe
\(f\in C^\infty([0,\infty);\mathcal S(\R^3;\R^3))\), and let
\((T_*,u,p)\) be the maximal pressure-complete classical history generated
by \((u_0,\nu,f)\). If \(T_*<\infty\), the whole-terminal record produced by
Theorem~\ref{thm:datum-intersection} projects onto
every finite \((N,M)\)-rectangle and gives
\begin{equation}
 \mathfrak R_{N,M}^*[u_0,\nu,f,T_*]
 :=\sup_{0<T<T_*}\mathfrak R_{N,M}(u;T)
 <\infty.
 \label{eq:datum-rectangle-bound}
\end{equation}
This assertion is coordinatewise in \((N,M)\); its finite constant may depend
on the selected temporal and spatial orders.  Equivalently, for every
\(n,m\in\Nzero\),
\begin{equation}
 V_n\in L^2(0,T_*;H^{m+1}),
 \qquad
 V_{n+1}\in L^2(0,T_*;H^{m-1}).
 \label{eq:datum-adjacent-rows}
\end{equation}
Every coordinate in \eqref{eq:datum-adjacent-rows} is the corresponding
distributional derivative of the single maximal history.

The independent rectangle relationship then estimates the complete equation.
For every \(n\in\Nzero\) and \(s\ge0\), put
\begin{align}
 \mathfrak D_{n,s}(S):={}&
 \nu\norm{\Lambda^{s+1}V_n(S)}_2^2
 +\int_0^S\!\left(
   \norm{\Lambda^sV_{n+1}}_2^2
   +\nu^2\norm{\Lambda^{s+2}V_n}_2^2
   +\norm{\Lambda^s\nabla Q_n}_2^2\right)\dd t.
 \label{eq:datum-complete-row-record}
\end{align}
Then
\begin{equation}
 \sup_{0<S<T_*}\mathfrak D_{n,s}(S)
 \le
 \nu\norm{\Lambda^{s+1}V_n^0}_2^2
 +\norm{\Lambda^s(F_n-B_n)}_{L^2(0,T_*;L^2)}^2
 <\infty.
 \label{eq:datum-complete-row-projection-bound}
\end{equation}
Thus the selected rectangle bounds, on the same terminal interval, its
adjacent temporal face, viscous face, full binomial nonlinear row, normalized
pressure-gradient face, and prescribed-force face.
\end{theorem}

\begin{proof}
Let \(I=(0,T_*)\).  Theorem~\ref{thm:datum-intersection} gives
\[
 V_n\in L^2(I;H^{M+1}),\qquad
 V_{n+1}\in L^2(I;H^{M-1}),
 \qquad 0\le n\le N.
\]
For \(M=0\), the second assertion is read through
\(V_{n+1}\in L^2(I;L^2)\hookrightarrow L^2(I;H^{-1})\).
The finitely many selected coordinates therefore have the finite total
\[
 C_{N,M}:=\sum_{n=0}^N
 \left(\norm{V_n}_{L^2(I;H^{M+1})}^2
       +\norm{V_{n+1}}_{L^2(I;H^{M-1})}^2\right)<\infty.
\]
Restriction decreases each norm, so
\(\mathfrak R_{N,M}(u;T)\le C_{N,M}\) for \(0<T<T_*\).
Taking the supremum proves \eqref{eq:datum-rectangle-bound}; monotone
exhaustion of the nonnegative norm integrals gives in fact
\(\mathfrak R_{N,M}^*=C_{N,M}\).  This projection uses the intersection
only to identify and place the coordinates; it does not derive the viscous
rectangle relationship.

Fix \((n,s)\).  Project the same intersection onto a finite rectangle high
enough to contain every velocity factor in \(B_n\) at spatial order \(s+2\).
The adjacent trace lemma gives the required \(L^\infty_tH^{s+2}_x\) factor,
while the higher spatial face gives the other factor in
\(L^2_tH^{s+2}_x\).  Sobolev multiplication, summed over the finite binomial
row, gives the explicit same-history estimate
\begin{equation}
 \norm{B_n}_{L^2(I;\dot H^s)}
 \le C_s\sum_{a=0}^{n}\binom na
 \norm{V_a}_{L^\infty(I;H^{s+2})}
 \norm{V_{n-a}}_{L^2(I;H^{s+2})}
 <\infty.
 \label{eq:datum-complete-binomial-projection-bound}
\end{equation}
The prescribed Schwartz force gives
\(F_n\in L^2(I;\dot H^s)\).  Now evaluate the independently derived
positive-sign identity \eqref{eq:complete-forced-rectangle-integrated} on
these selected coordinates.  Its right-hand side is finite, and taking the
supremum over \(S<T_*\) proves
\eqref{eq:datum-complete-row-projection-bound}.  The pressure term is present
because \eqref{eq:complete-forced-pressure-square} is the unprojected
orthogonal completion of the same row.  No term of the forced equation and no
coordinate of another history has entered the estimate.
\end{proof}

\begin{proposition}[Complete momentum rows inside one finite rectangle]
\label{prop:rectangle-bounds-complete-rows}
Assume the hypotheses of Theorem~\ref{thm:datum-rectangles}.  Fix
\(N\in\Nzero\), \(M\ge2\), and \(0<T\le T_*\), and put
\begin{align*}
 R_v&:=\sum_{n=0}^N\left(
 \norm{V_n}_{L^2(0,T;H^{M+1})}^2
 +\norm{V_{n+1}}_{L^2(0,T;H^{M-1})}^2\right),\\
 I_{N,M}&:=\sum_{n=0}^N\norm{V_n^0}_{H^M}^2.
\end{align*}
Then the adjacent faces of this same rectangle give
\begin{equation}
 \sum_{n=0}^N\sup_{0\le t\le T}\norm{V_n(t)}_{H^M}^2
 \le I_{N,M}+R_v,
 \label{eq:rectangle-adjacent-trace-bound}
\end{equation}
and the complete binomial nonlinear rows obey
\begin{equation}
 \sum_{n=0}^N\norm{B_n}_{L^2(0,T;H^{M-1})}^2
 \le C_{N,M}(I_{N,M}+R_v)R_v.
 \label{eq:rectangle-complete-binomial-bound}
\end{equation}
For \(0<S\le T\), define the pressure-complete dissipation record
\begin{align*}
 \mathcal D_{N,M}(S):=\sum_{n=0}^N\bigg[&
 \nu\norm{V_n(S)}_{\dot H^M}^2
 +\norm{V_{n+1}}_{L^2(0,S;\dot H^{M-1})}^2\\
 &+\nu^2\norm{V_n}_{L^2(0,S;\dot H^{M+1})}^2
 +\norm{\nabla Q_n}_{L^2(0,S;\dot H^{M-1})}^2\bigg].
\end{align*}
The integrated complete rows satisfy
\begin{align}
 \mathcal D_{N,M}(S)
 \le{}&\nu I_{N,M}
 +2\sum_{n=0}^N\norm{F_n}_{L^2(0,T;H^{M-1})}^2
 \notag\\
 &+2C_{N,M}(I_{N,M}+R_v)R_v.
 \label{eq:rectangle-complete-row-bound}
\end{align}
Thus one finite velocity rectangle, together with the displayed prescribed
force coordinates, bounds on the same interval and in the same equations its
adjacent temporal rows, viscous Laplacians, complete nonlinear rows, and
normalized pressure gradients.  Lower spatial orders follow by Sobolev nesting
from a rectangle with \(M\ge2\).
\end{proposition}

\begin{proof}
For \(0\le n\le N\), the Hilbert triple
\(H^{M+1}\subset H^M\subset H^{M-1}\) and
\(\partial_tV_n=V_{n+1}\) give
\[
 \norm{V_n(t)}_{H^M}^2-\norm{V_n^0}_{H^M}^2
 =2\operatorname{Re}\int_0^t
 \ip{(1-\Delta)^{(M+1)/2}V_n}
 {(1-\Delta)^{(M-1)/2}V_{n+1}}_{L^2}\dd s.
\]
Cauchy--Schwarz in time and \(2ab\le a^2+b^2\), followed by summation in
\(n\), prove \eqref{eq:rectangle-adjacent-trace-bound}.  Sobolev
multiplication applied to every allocation in \(B_n\), using one factor in
\(L^\infty(0,T;H^M)\) and the other in
\(L^2(0,T;H^{M+1})\), proves
\eqref{eq:rectangle-complete-binomial-bound}.  Finally, sum
\eqref{eq:complete-forced-rectangle-integrated} over \(0\le n\le N\) at
\(s=M-1\), use
\(\norm{F_n-B_n}^2\le2\norm{F_n}^2+2\norm{B_n}^2\), and substitute
\eqref{eq:rectangle-complete-binomial-bound}.  This gives
\eqref{eq:rectangle-complete-row-bound}.  The identities
\[
 \nabla Q_n=(I-\Pp)(F_n-B_n),
 \qquad
 V_{n+1}=\nu\Delta V_n-\Pp B_n+\Pp F_n
\]
show directly that these are the terms of the same pressure-complete row.
\end{proof}

The finite sum above evaluates the coordinates selected by
Definition~\ref{def:finite-rectangle}.  Its finiteness comes from the
datum-generated intersection, while the complete-equation estimate retains
and relates all faces of those coordinates.  The next proposition records the
exact monotone-exhaustion identity and the equivalent memberships.

\Needspace{14\baselineskip}
\begin{proposition}[Norm equivalence and whole-terminal realization]
\label{prop:datum-terminal-realization}

For the fixed maximal history, the family of bounds
\eqref{eq:datum-rectangle-bound} is equivalent to the adjacent-row statement
\eqref{eq:datum-adjacent-rows}.  Whenever either formulation holds, for every
\(N,M\in\Nzero\),
\begin{equation}
 \sum_{n=0}^{N}\left(
 \norm{V_n}_{L^2(0,T_*;H^{M+1})}^2
 +\norm{V_{n+1}}_{L^2(0,T_*;H^{M-1})}^2
 \right)
 =\lim_{T\uparrow T_*}\mathfrak R_{N,M}(u;T)
 =\mathfrak R_{N,M}^*[u_0,\nu,f,T_*].
 \label{eq:datum-terminal-monotone-realization}
\end{equation}
Every row is the corresponding distributional derivative of the single
pressure-normalized maximal history.
\end{proposition}

\begin{proof}
Assume first \eqref{eq:datum-adjacent-rows}.  For fixed \(N,M\), set
\[
 C_{N,M}:=
 \sum_{n=0}^{N}\left(
  \norm{V_n}_{L^2(0,T_*;H^{M+1})}^2
  +\norm{V_{n+1}}_{L^2(0,T_*;H^{M-1})}^2
 \right).
\]
The sum is finite because it contains finitely many instances of
\eqref{eq:datum-adjacent-rows}.  Restriction from \((0,T_*)\) to \((0,T)\)
gives
\[
 \mathfrak R_{N,M}(u;T)\le C_{N,M}
 \qquad(0<T<T_*),
\]
and hence \eqref{eq:datum-rectangle-bound}.

Conversely, assume \eqref{eq:datum-rectangle-bound}.  Every term in
\(\mathfrak R_{N,M}(u;T)\) is the integral of a nonnegative function.
Monotone exhaustion of \((0,T_*)\) gives
\eqref{eq:datum-terminal-monotone-realization}.  Taking \(N=n\) and \(M=m\)
yields \eqref{eq:datum-adjacent-rows} for each finite \((n,m)\).
Proposition~\ref{prop:formal-mixed-recurrences} identifies every entry with
the corresponding distributional derivative of the fixed pressure-normalized
history.
\end{proof}

\begin{corollary}[Whole-terminal completion of every critical height]
\label{cor:datum-critical-height-completion}
Under the hypotheses of Theorem~\ref{thm:datum-rectangles}, every completed
critical coordinate belongs to \(L^1(0,T_*)\), and
\begin{equation}
 \mathscr H_k[u_0,\nu,f,T_*]
 :=\norm{R_k}_{L^1(0,T_*)}
 =\frac12\norm{\Lambda^{k+1/2}u}_{L^2(0,T_*;L^2)}^2
 \le\frac12\mathfrak R_{0,k}^*[u_0,\nu,f,T_*]
 \qquad(k\in\Nzero).
 \label{eq:datum-completed-height-bound}
\end{equation}
For every \(m\in\Nzero\), the spatial component of
\eqref{eq:independent-mixed-intersection} satisfies the completed-height estimate
\begin{equation}
 \norm u_{L^2(0,T_*;H^m)}^2
 \le
 2^m\left(
  T_*K_f(T_*)^2+\mathfrak R_{0,m}^*[u_0,\nu,f,T_*]
 \right).
 \label{eq:datum-height-to-spatial-column}
\end{equation}
Thus
\begin{equation}
 (R_k)_{k\in\Nzero}\in\prod_{k\in\Nzero}L^1(0,T_*),
 \qquad
 u\in\bigcap_{m\in\Nzero}L^2(0,T_*;H^m(\R^3)).
 \label{eq:datum-complete-height-and-spatial-columns}
\end{equation}
\end{corollary}

\begin{samepage}
\begin{proof}
Take \(N=0\) and \(M=k\) in
\eqref{eq:datum-terminal-monotone-realization}.  Since
\(\abs\xi^{2k+1}\le\langle\xi\rangle^{2k+2}\),
\begin{align*}
 2\norm{R_k}_{L^1(0,T_*)}
 &=\norm{\Lambda^{k+1/2}u}_{L^2(0,T_*;L^2)}^2\\
 &\le\norm u_{L^2(0,T_*;H^{k+1})}^2
 \le\mathfrak R_{0,k}^*[u_0,\nu,f,T_*],
\end{align*}
which proves \eqref{eq:datum-completed-height-bound}.  For
\(\abs\xi\le1\),
\(\langle\xi\rangle^{2m}\le2^m\); for \(\abs\xi>1\),
\(\langle\xi\rangle^{2m}\le2^m\abs\xi^{2m+1}\).  Hence
\[
 \norm v_{H^m}^2
 \le2^m\left(\norm v_2^2+
 \norm{\Lambda^{m+1/2}v}_2^2\right).
\]
Integrating this inequality, using
\eqref{eq:energy-Linf} and
\eqref{eq:datum-completed-height-bound}, proves
\eqref{eq:datum-height-to-spatial-column}.  Since \(k\) and \(m\) are
arbitrary finite indices, \eqref{eq:datum-complete-height-and-spatial-columns}
follows.
\end{proof}
\end{samepage}

\begin{proposition}[Pressure completion]
\label{prop:rectangle-pressure-completion}
Under the hypotheses of Theorem~\ref{thm:datum-rectangles}, for every
\(N,M\in\Nzero\),
\begin{equation}
 \mathfrak Q_{N,M}^*[u_0,\nu,f,T_*]
 :=\sup_{0<T<T_*}\mathfrak Q_{N,M}(p;T)
 \le
 C_M^{\rm p}(2^{N+1}-1)\mathfrak R_{N,M}^*[u_0,\nu,f,T_*]
 +\sum_{n=0}^N\mathfrak F^{\rm p}_{n,0,M}(0,T_*).
 \label{eq:whole-terminal-pressure-bound}
\end{equation}
\begin{samepage}
For every \(0\le n\le N\), the identities
\begin{align}
 Q_n
 &=R_iR_j\sum_{a=0}^n\binom naV_{a,i}V_{n-a,j}
 -(-\Delta)^{-1}\nabla\cdot F_n,
 \label{eq:rectangle-pressure-identity}\\
 V_{n+1}
 &=\nu\Delta V_n
 -\sum_{a=0}^n\binom na(V_a\cdot\nabla)V_{n-a}-\nabla Q_n+F_n
 \label{eq:rectangle-velocity-identity}
\end{align}
hold respectively in
\(L^1(0,T_*;H^M)\) and \(L^1(0,T_*;H^{M-1})\), including when
\(M=0\).
\end{samepage}
\end{proposition}

\begin{proof}
For every \(0<T<T_*\),
Theorem~\ref{thm:formal-finite-row-pressure} gives
\[
 \norm{Q_n}_{L^1(0,T;H^M)}
 +\norm{\nabla Q_n}_{L^1(0,T;H^{M-1})}
 \le C_M^{\rm p}2^n\mathfrak R_{N,M}^*[u_0,\nu,f,T_*]
 +\mathfrak F^{\rm p}_{n,0,M}(0,T).
\]
Sum over \(0\le n\le N\), then let \(T\uparrow T_*\).  Monotone
convergence in the nonnegative norm integrals proves
\eqref{eq:whole-terminal-pressure-bound}.  Since
\(\nabla\cdot V_a=0\),
\[
 (V_a\cdot\nabla)V_{n-a}
 =\nabla\cdot(V_{n-a}\otimes V_a).
\]
Lemma~\ref{lem:sobolev-product} and Cauchy--Schwarz in time give
\[
 \norm{V_{n-a}\otimes V_a}_{L^1(0,T_*;H^M)}
 \le C_M^{\rm alg}
 \norm{V_a}_{L^2(0,T_*;H^{M+1})}
 \norm{V_{n-a}}_{L^2(0,T_*;H^{M+1})}.
\]
Thus every transport term belongs to
\(L^1(0,T_*;H^{M-1})\), also for \(M=0\).  Moreover,
\[
 \norm{\Delta V_n}_{L^1H^{M-1}}
 \le T_*^{1/2}\norm{V_n}_{L^2H^{M+1}}.
\]
Theorem~\ref{thm:datum-rectangles} places \(V_{n+1}\) in
\(L^2(0,T_*;H^{M-1})\), hence in \(L^1(0,T_*;H^{M-1})\).
The prescribed row \(F_n\) belongs to \(L^1(0,T_*;H^{M-1})\), and its
pressure potential belongs to the stated spaces by
\eqref{eq:force-pressure-potential-bound}.
The strict-interval identities are identities of distributions; their
whole-terminal representatives belong to the stated spaces.  This proves the
two identities on \((0,T_*)\).
\end{proof}

\begin{definition}[Finite-window spaces]
\label{def:finite-window-space}
For \(N,M\in\Nzero\), set
\[
 \mathscr X_{N,M}:=
 \prod_{n=0}^{N}L^2(0,T_*;H^{M+1})
 \times
 \prod_{n=0}^{N}L^2(0,T_*;H^{M-1})
 \times
 \prod_{n=0}^{N}L^1(0,T_*;H^M)
 \times
 \prod_{n=0}^{N}L^1(0,T_*;H^{M-1}).
\]
The coordinate projection \(\pi_{N,M}\) selects the rows through
\(N\), their next time derivatives, and the corresponding pressure and force
rows from \(\cI[u_0,\nu,f]\), in the displayed Sobolev spaces.
The whole-terminal rectangle is the element
\begin{equation}
 \cR_{N,M}[u_0,\nu,f;T_*]
 :=\pi_{N,M}\cI[u_0,\nu,f]=\left(
 (V_n)_{n=0}^{N},
 (V_{n+1})_{n=0}^{N};
 (Q_n)_{n=0}^{N};
 (F_n)_{n=0}^{N}
 \right)\in\mathscr X_{N,M}.
 \label{eq:whole-terminal-rectangle-element}
\end{equation}
If \(N'\le N\) and \(M'\le M\), let
\[
 \rho_{N',M'}^{N,M}:\mathscr X_{N,M}\longrightarrow\mathscr X_{N',M'}
\]
discard the temporal coordinates with index larger than \(N'\) and apply
\[
 H^{M+1}\hookrightarrow H^{M'+1},
 \qquad H^{M-1}\hookrightarrow H^{M'-1},
 \qquad H^M\hookrightarrow H^{M'},
 \qquad H^{M-1}\hookrightarrow H^{M'-1}
\]
to the retained coordinates.
\end{definition}

\begin{proposition}[Restriction compatibility]
\label{prop:rectangle-compatibility}
For \(N'\le N\) and \(M'\le M\),
\begin{equation}
 \rho_{N',M'}^{N,M}\cR_{N,M}[u_0,\nu,f;T_*]
 =\cR_{N',M'}[u_0,\nu,f;T_*].
 \label{eq:rectangle-compatibility}
\end{equation}
If
\((N'',M'')\le(N',M')\le(N,M)\), then
\[
 \rho_{N'',M''}^{N',M'}\rho_{N',M'}^{N,M}
 =\rho_{N'',M''}^{N,M}.
\]
Thus the spaces \(\mathscr X_{N,M}\) and the maps
\(\rho_{N',M'}^{N,M}\) form an inverse system, and the datum-generated
rectangles form a compatible family in it.
\end{proposition}

\begin{proof}
All retained entries are the fixed distributions
\(\partial_t^nu\), \(\partial_t^{n+1}u\), \(\partial_t^np\), and
\(\partial_t^nf\).
The Sobolev embeddings preserve those distributions.  Deletion of rows and
composition of canonical embeddings are associative.
\end{proof}
